% Figure: superposition of axial and bending stress through the depth of a
% section, and the neutral-axis shift the axial part produces. Three
% through-depth stress profiles plotted with depth y on the vertical axis
% and stress on the horizontal: the uniform axial part (constant), the
% linear bending part (antisymmetric about the centroid), and their sum.
% The zero crossing of the sum is off the centroid: that is the shifted
% neutral axis. Numbers follow the opener worked example: axial 37.5 MPa,
% bending +/-187 MPa at the extreme fibres, sum 224 (tension) and -150
% (compression). Depth y runs -1 (top, compression fibre) to +1 (bottom).
\begin{figure}[htbp]
\centering
\begin{tikzpicture}
\begin{axis}[
  width=13cm, height=7.5cm,
  xlabel={stress $\sigma$ (MPa)},
  ylabel={position through depth, $y/(h/2)$},
  xmin=-220, xmax=260, ymin=-1.15, ymax=1.15,
  axis lines=middle,
  axis line style={->},
  tick label style={font=\scriptsize},
  label style={font=\small},
  legend pos=outer north east,
  legend style={font=\scriptsize},
]
% axial: constant 37.5 MPa (vertical line in this plot: sigma fixed, y free)
\addplot[englime!60!black, very thick] coordinates {(37.5,-1) (37.5,1)};
\addlegendentry{axial $\sigma_N = 37.5$}
% bending: sigma = -187 y  (top y=+1 -> compression -187 on top? use
% convention: bottom fibre y=+1 tension +187, top y=-1 compression -187)
\addplot[engblue, very thick] coordinates {(-187,-1) (187,1)};
\addlegendentry{bending $\sigma_M = 187\,y$}
% sum: 37.5 + 187 y : at y=-1 -> -149.5, at y=+1 -> 224.5
\addplot[engred, very thick] coordinates {(-149.5,-1) (224.5,1)};
\addlegendentry{sum $\sigma_N + \sigma_M$}
% mark the shifted neutral axis: sum=0 -> y = -37.5/187 = -0.20
\filldraw[engorange] (axis cs:0,-0.2005) circle[radius=2pt];
\node[engorange, font=\scriptsize, anchor=west] at (axis cs:8,-0.31)
  {shifted neutral axis, $y=-0.20$};
% centroid marker
\filldraw[enggray] (axis cs:0,0) circle[radius=1.4pt];
\node[enggray, font=\scriptsize, anchor=south east] at (axis cs:-4,0.02) {centroid};
% extreme-fibre value labels
\node[engred, font=\scriptsize, anchor=south] at (axis cs:224.5,1.0) {$224$};
\node[engred, font=\scriptsize, anchor=north] at (axis cs:-149.5,-1.0) {$-150$};
\end{axis}
\end{tikzpicture}
\caption[Superposition of axial and bending stress]{Stress through the
depth of the opener's cantilever section, plotted with depth on the
vertical axis and stress on the horizontal. The axial pull contributes a
constant $37.5\ \text{MPa}$ (vertical green line); the bending moment
contributes a linear profile antisymmetric about the centroid, $\pm 187\
\text{MPa}$ at the extreme fibres; their sum is the tilted red line. The
constant axial part slides the whole bending line sideways, so the sum
crosses zero not at the centroid but at $y = -0.20$ of the half-depth: the
neutral axis has moved off the centroid. The extreme-fibre stresses are
$224\ \text{MPa}$ in tension on the bottom and $150\ \text{MPa}$ in
compression on the top, the three-formula superposition that is the daily
arithmetic of combined loading.}
\label{fig:vol03ch03:combined-axial-bending-profile}
\end{figure}
